\begin{titlepage}
    \centering
    \vspace*{1cm}
        
    \Huge
    \textbf{Longitudinal Space-Charge Effects of Transverse Motion in the PS}
        
    \vspace{0.5cm}
    \LARGE
    % subtitle
        
    \vspace{1.5cm}
        
    \textbf{Alexander J. Laut}
        
    \vfill
        
    \includegraphics[width=0.4\textwidth]{logos/LogoOutline-Bleu.pdf}
    \vfill
    
    A thesis presented for the degree of\\
    Masters of Physics
    
    \vspace{0.8cm}
        
    \includegraphics[width=0.4\textwidth]{logos/logo.png}
    
    \Large
    Department of Physics\\
    % University Paris-Saclay\\
    France\\
    September 2021
    
    % Composition of the jury:
    % \begin{itemize}
    %     \item Nicolas Chauvin
    %     \item Nicolas Delerue,
    %     \item Sophi Kazamias IJC Lab
    % \end{itemize}
\end{titlepage}

\begin{abstract}
Particles in an intense bunch will experience longitudinal self-fields due to space-charge. This space charge effect, described by geometric factors dependent on a particle’s transverse position, beam size, and beam pipe aperture, is usually incorporated into longitudinal trackers on a per-turn basis by assuming particles are uniformly affected. Sampling from a 6D distribution in phase-space, a particle's transverse trajectories within a bunch will vary due to Betatron motion. A transverse tracker was developed to characterize the effective geometric factor of a given particle accounting for its transverse emittance, phase advance, and dispersion within a ring. A particle’s effective geometry factor is then estimated by interpolation without the need for transverse tracking. This variance in geometric factor due to transverse motion was incorporated into the longitudinal tracker BLonD, systematically affecting the synchrotron frequency distribution. The filamentation rate of a mismatched or perturbed longitudinal distribution was increased, which could suggest a stabilizing space-charge phenomenon.
\end{abstract}

\chapter*{Acknowledgements}
\thispagestyle{empty}

I'd like to thank my supervisor Dr. Alexandre Lasheen for providing me this amazing opportunity to work at CERN and be extraordinarily accommodating to the many challenges faced during the COVID-19 pandemic. He provided substantial guidance throughout understanding the topic, and concisely acted as an excellent resource.